\section{\texorpdfstring{Cartan–\allowbreak{}Eilenberg resolutions,\allowbreak{} composition,\allowbreak{} and resolution functors}{   and resolution functors}}\label{proofs-7096-7578-cartaneilenberg-resolutions-composition-and-resolution-functors}

This is editorial review material,\allowbreak{} not a replacement exposition for the translated source. Source:\allowbreak{} the Stacks Project authors,\allowbreak{} Derived Categories,\allowbreak{} original commit a04446e5\allowbreak{}7ec1fbc2\allowbreak{}52a871af\allowbreak{}cec7752f\allowbreak{}b2807b14, lines 7096–\allowbreak{}7578. The frozen original and the current source have the identities recorded in SOURCE\_\allowbreak{}READING\_\allowbreak{}PART10.json. Original line numbers below refer to that original. The current Homology chapter supplies the commuting-square double-complex convention at 4207–\allowbreak{}4223 and the total differential at 4283–\allowbreak{}4300. Complete earlier arguments used here are retained in PROOFS\_\allowbreak{}6209\_\allowbreak{}7089.md,\allowbreak{} especially sections 2–\allowbreak{}7,\allowbreak{} and PROOFS\_\allowbreak{}5091\_\allowbreak{}5900.md with ADDENDUM\_\allowbreak{}01. No novelty is claimed.

\subsection{\texorpdfstring{1. Existence,\allowbreak{} the actual double-complex maps,\allowbreak{} and prescribed resolutions}{1.  the actual double-complex  and prescribed resolutions}}\label{proofs-7096-7578-existence-the-actual-double-complex-maps-and-prescribed-resolutions}

Keep the source\textquotesingle s complex K,\allowbreak{} with K\textasciicircum{}p\allowbreak{}=0 for p<n,\allowbreak{} and its actual objects Z\textasciicircum{}p\allowbreak{}=ker(d\_\allowbreak{}K\textasciicircum{}p)\allowbreak{},\allowbreak{} B\textasciicircum{}p\allowbreak{}=im(d\_\allowbreak{}K\textasciicircum{}\{p-1\})\allowbreak{},\allowbreak{} H\textasciicircum{}p\allowbreak{}=Z\textasciicircum{}p/\allowbreak{}B\textasciicircum{}p. Its two short exact sequences,\allowbreak{} in every degree p>\allowbreak{}=n,\allowbreak{} are

\[
 0\longrightarrow B^p\longrightarrow Z^p\longrightarrow H^p
       \longrightarrow0,\qquad
 0\longrightarrow Z^p\longrightarrow K^p\longrightarrow B^{p+1}
       \longrightarrow0.                                      \tag{1.1}
\]

Here B\textasciicircum{}n\allowbreak{}=0 and Z\textasciicircum{}n\allowbreak{}=H\textasciicircum{}n. The source starts by resolving Z\^{}n and applies the injective horseshoe construction successively along (1.1)\allowbreak{},\allowbreak{} retaining the first resolution at each step. This is the already proved construction in section 6 of PROOFS\_\allowbreak{}6209\_\allowbreak{}7089.md. It gives exact sequences of complexes

\[
 0\to J_B^{p,\bullet}\xrightarrow{b_p}J_Z^{p,\bullet}
       \xrightarrow{c_p}J_H^{p,\bullet}\to0,\qquad
 0\to J_Z^{p,\bullet}\xrightarrow{z_p}I^{p,\bullet}
       \xrightarrow{e_p}J_B^{p+1,\bullet}\to0.                   \tag{1.2}
\]

All the complexes in (1.2)\allowbreak{} have injective terms and vanish in vertical degrees q\textless0. Their augmentations form morphisms from the two source sequences. Put

\[
 d_1^{p,\bullet}=z_{p+1}b_{p+1}e_p,\qquad
 d_2^{p,q}=d_{I^{p,\bullet}}^q.                                \tag{1.3}
\]

These are the source\textquotesingle s maps,\allowbreak{} with their full factorization. Since e\_\allowbreak{}\{p\allowbreak{}+1\}z\_\allowbreak{}\{p\allowbreak{}+1\}\allowbreak{}=0,\allowbreak{} consecutive d\_\allowbreak{}1's compose to zero. Every factor in (1.3)\allowbreak{} is a vertical chain map,\allowbreak{} so d\_\allowbreak{}1d\_\allowbreak{}2\allowbreak{}=d\_\allowbreak{}2d\_\allowbreak{}1,\allowbreak{} with no inserted minus sign. Since e\_\allowbreak{}p is epic and z\_\allowbreak{}\{p\allowbreak{}+1\}b\_\allowbreak{}\{p\allowbreak{}+1\} is monic,\allowbreak{} its kernel is J\_\allowbreak{}Z\textasciicircum{}\{p,\allowbreak{}\textbackslash{}bullet\} and its image is J\_\allowbreak{}B\textasciicircum{}\{p\allowbreak{}+1,\allowbreak{}\textbackslash{}bullet\},\allowbreak{} with the actual inclusions from (1.2)\allowbreak{}. The horizontal cohomology in degree p is J\_\allowbreak{}Z\textasciicircum{}\{p,\allowbreak{}\textbackslash{}bullet\}/\allowbreak{}J\_\allowbreak{}B\textasciicircum{}\{p,\allowbreak{}\textbackslash{}bullet\}\allowbreak{}=J\_\allowbreak{}H\textasciicircum{}\{p,\allowbreak{}\textbackslash{}bullet\}. These identifications commute with d\_\allowbreak{}2. The augmentations therefore give every resolution required in original 7116–\allowbreak{}7122,\allowbreak{} including the image resolution indexed by p\allowbreak{}+1 rather than p in the B notation. Augmentation compatibility in (1.2)\allowbreak{} gives d\_\allowbreak{}1 epsilon\allowbreak{}=epsilon d\_\allowbreak{}K. Set every column with p\textless n equal to zero. This proves the stated existence theorem.

\textbf{DERIVED-UNDER-012:\allowbreak{} the boundary and cohomology resolutions may be prescribed simultaneously.} Fix injective resolutions J\_\allowbreak{}B\textasciicircum{}\{p,\allowbreak{}\textbackslash{}bullet\} of B\^{}p for every p>\allowbreak{}=n\allowbreak{}+1 and J\_\allowbreak{}H\textasciicircum{}\{p,\allowbreak{}\textbackslash{}bullet\} of H\^{}p for every p>\allowbreak{}=n,\allowbreak{} all vanishing for q\textless0. Set J\_\allowbreak{}B\textasciicircum{}\{p,\allowbreak{}\textbackslash{}bullet\}\allowbreak{}=0 for p<\allowbreak{}=n and J\_\allowbreak{}H\textasciicircum{}\{p,\allowbreak{}\textbackslash{}bullet\}\allowbreak{}=0 for p\textless n. Apply the both-prescribed-outer-resolution construction DERIVED-UNDER-010 first to the first sequence of (1.1)\allowbreak{},\allowbreak{} using the two given outer resolutions,\allowbreak{} to construct J\_\allowbreak{}Z\textasciicircum{}\{p,\allowbreak{}\textbackslash{}bullet\}. Apply it to the second sequence using that J\_\allowbreak{}Z\textasciicircum{}\{p,\allowbreak{}\textbackslash{}bullet\} and the already given J\_\allowbreak{}B\textasciicircum{}\{p\allowbreak{}+1,\allowbreak{}\textbackslash{}bullet\},\allowbreak{} to construct I\textasciicircum{}\{p,\allowbreak{}\textbackslash{}bullet\}. Thus both sequences (1.2)\allowbreak{} exist for each p with exactly the given outer complexes and augmentations. The proof of (1.3)\allowbreak{} then applies verbatim. In particular the resulting horizontal boundary and cohomology resolution complexes identify with the prescribed ones,\allowbreak{} not merely with unspecified quasi-isomorphic complexes.

The earlier horseshoe proof constructs a middle differential \[
 \begin{pmatrix}d_{J_1}&\delta'\\0&d_{J_3}\end{pmatrix}
\] and an augmentation through the actual matrix \(\left(\begin{smallmatrix}1&h\\0&f_3\end{smallmatrix}\right)\). The relation \(\delta'f_3-\delta=-d_{J_1}h+h\,d_C\) proves that augmentation is a chain map. We use that construction,\allowbreak{} including h; an arbitrary pair of chosen resolutions is not treated as if its connecting maps already agreed strictly. Both outer resolutions are zero below 0,\allowbreak{} so the constructed middle complexes are also zero below 0. Each step uses finite direct sums only. Once the prescribed resolutions exist,\allowbreak{} no further global enough-injectives hypothesis is needed for this version.

There is a necessary boundary restriction:\allowbreak{} a nonzero acyclic injective resolution of B\textasciicircum{}n\allowbreak{}=0 cannot be prescribed while requiring all columns p\textless n to be zero. The image of d\_\allowbreak{}1\textasciicircum{}\{n-1,\allowbreak{}\textbackslash{}bullet\} is then literally zero. Our statement explicitly sets that boundary complex to zero. It does not assert that arbitrary resolutions at that boundary work.

DERIVED-RECON-043 concerns only the article at 7113. Both reports P02-E0048 and DERIVED-040 describe the existing one-operation MC-STK-ERR-0851 correction from ``a i'' to ``an i''. The lower-bound condition itself is unchanged and valid.

\subsection{\texorpdfstring{2. Totalization,\allowbreak{} both spectral sequences,\allowbreak{} and the reported index error}{2.  both spectral  and the reported index error}}\label{proofs-7096-7578-totalization-both-spectral-sequences-and-the-reported-index-error}

Write T\allowbreak{}=Tot(I)\allowbreak{}. In total degree m the only possibly nonzero summands are I\textasciicircum{}\{p,\allowbreak{}m-p\} for n<\allowbreak{}=p<\allowbreak{}=m. Thus T is zero below n and has injective terms:\allowbreak{} a finite direct sum of injectives is injective,\allowbreak{} since extension of a map into it is extension of each of finitely many components. Its differential on I\textasciicircum{}\{p,\allowbreak{}q\} is exactly

\[
 d_T=d_1^{p,q}+(-1)^p d_2^{p,q}.                               \tag{2.1}
\]

The vertical augmentation equation d\_\allowbreak{}2\textasciicircum{}\{p,\allowbreak{}0\}epsilon\textasciicircum{}p\allowbreak{}=0 and the horizontal equation from section 1 prove that epsilon:\allowbreak{}K\allowbreak{}->T is a chain map. The first spectral sequence for I has E\_\allowbreak{}1\textasciicircum{}\{p,\allowbreak{}0\}\allowbreak{}=K\textasciicircum{}p,\allowbreak{} E\_\allowbreak{}1\textasciicircum{}\{p,\allowbreak{}q\}\allowbreak{}=0 for q!\allowbreak{}=0,\allowbreak{} and d\_\allowbreak{}1\allowbreak{}=d\_\allowbreak{}K. Finite diagonal support gives convergence. Consequently its E\_\allowbreak{}2 and E\_\allowbreak{}infinity terms are H\textasciicircum{}p(K)\allowbreak{} on q\allowbreak{}=0 and zero elsewhere. Apply the same construction to the double complex having K in row 0 and zero other rows. The augmentation is a map of double complexes and induces the identity under these E\_\allowbreak{}1 identifications. The associated cohomology filtrations have only one nonzero quotient in each total degree. Hence H\textasciicircum{}m(epsilon)\allowbreak{} itself is the resulting isomorphism. This verifies the map omitted at current Homology 6833–\allowbreak{}6835,\allowbreak{} rather than just an abstract isomorphism of objects. Therefore T is an injective resolution of K.

Let F:\allowbreak{}A\allowbreak{}->B be an additive functor of abelian categories. For the source theorem F is also left exact; we first keep that assumption and then track where it is used. The individual injective resolutions already constructed define RF at K\textasciicircum{}p[0]\allowbreak{},\allowbreak{} H\textasciicircum{}p(K)\allowbreak{}[0]\allowbreak{},\allowbreak{} and K,\allowbreak{} by the injective terminal-comparison argument proved in section 7 of the preceding note. No assumption that RF exists on every other object is needed. Finite additivity gives the canonical degreewise isomorphism F(Tot I)\allowbreak{}\allowbreak{}->Tot(F(I)\allowbreak{})\allowbreak{},\allowbreak{} preserving (2.1)\allowbreak{} and its sign. Put L\allowbreak{}=F(I)\allowbreak{}.

For the first filtration,\allowbreak{} current Homology 6704–\allowbreak{}6709 gives

\[
 {}'E_0^{p,q}=F(I^{p,q}),\quad {}'d_0=(-1)^p F(d_2),\quad
 {}'E_1^{p,q}=H^q(F(I^{p,\bullet}))=R^qF(K^p).                  \tag{2.2}
\]

Multiplication of a differential by -1 leaves its kernel and image unchanged; this identifies the displayed cohomology without losing the sign in the differential. Its d\_\allowbreak{}1 is the map R\textasciicircum{}qF(d\_\allowbreak{}K\textasciicircum{}p)\allowbreak{}. Indeed d\_\allowbreak{}1\textasciicircum{}\{p,\allowbreak{}\textbackslash{}bullet\} is a comparison of the injective resolutions of K\^{}p and K\textasciicircum{}\{p\allowbreak{}+1\} extending d\_\allowbreak{}K\textasciicircum{}p; the uniqueness on homotopy classes proved earlier identifies its induced derived map.

Fix a vertical degree p. In horizontal degree q the two exact sequences of objects are

\[
 0\to B_I^{q,p}\to Z_I^{q,p}\to H_I^{q,p}\to0,\qquad
 0\to Z_I^{q,p}\to I^{q,p}\to B_I^{q+1,p}\to0.                 \tag{2.3}
\]

They split because their first terms are injective. Applying F preserves those split sequences. Thus the kernel,\allowbreak{} image,\allowbreak{} and cohomology of F(d\_\allowbreak{}1)\allowbreak{} are respectively F(Z\_\allowbreak{}I)\allowbreak{},\allowbreak{} F(B\_\allowbreak{}I)\allowbreak{},\allowbreak{} and F(H\_\allowbreak{}I)\allowbreak{},\allowbreak{} via their actual induced inclusion and quotient maps. The cohomology identification is canonical:\allowbreak{} it is the map induced by F of the cycle inclusion and the quotient,\allowbreak{} whose kernel and image have just been proved; it does not depend on chosen splittings. Vertical chain maps preserve those inclusions and quotients. Hence

\[
 {}''E_1^{p,q}=F(H_I^q(I^{\bullet,p})),\qquad
 {}''d_1^{p,q}=(-1)^q F(d_{\,H_I^q(I^{\bullet,\bullet})}^{p}),
 \qquad
 {}''E_2^{p,q}=R^pF(H^q(K)).                                  \tag{2.4}
\]

This proves DERIVED-RECON-044 /\allowbreak{} P02-E0049. The two p\textquotesingle s at original 7248–\allowbreak{}7249 must be q:\allowbreak{} the resolution complex is H\_\allowbreak{}I\textasciicircum{}q(I)\allowbreak{},\allowbreak{} and the object it resolves is H\textasciicircum{}q(K)\allowbreak{}. The vertical cohomology degree p,\allowbreak{} the superscript p on E\_\allowbreak{}1,\allowbreak{} and the sign (-1)\allowbreak{}\textasciicircum{}q remain unchanged. The report\textquotesingle s phrase ``vertical homology degree'' and its final formula do not precisely describe both printed occurrences; the source and (2.4)\allowbreak{},\allowbreak{} not that wording,\allowbreak{} determine the two edits. No previous canonical correction covers either occurrence.

Both filtrations are finite in each degree and both spectral sequences converge to H\textasciicircum{}m(RF(K)\allowbreak{})\allowbreak{}. For a chosen column lower bound n,\allowbreak{} their supports are respectively p>\allowbreak{}=n,\allowbreak{}q>\allowbreak{}=0 and p>\allowbreak{}=0,\allowbreak{}q>\allowbreak{}=n. The differential has bidegree (r,\allowbreak{}1-r)\allowbreak{}. At a first-sequence term (p,\allowbreak{}q)\allowbreak{},\allowbreak{} an incoming differential is zero when r\textgreater p-n and an outgoing one is zero when r>q\allowbreak{}+1. At a second-sequence term they are zero when r\textgreater p and r>q-n\allowbreak{}+1 respectively. On total degree m>\allowbreak{}=n,\allowbreak{} every term has therefore stabilized by page r\allowbreak{}=m-n\allowbreak{}+2. Each resulting filtration has at most m-n\allowbreak{}+1 potentially nonzero graded pieces. All terms and the abutment vanish in total degrees m\textless n.

\textbf{DERIVED-UNDER-011,\allowbreak{} first part.} All arguments in this section used additivity,\allowbreak{} the given injective resolutions,\allowbreak{} and the split sequences (2.3)\allowbreak{},\allowbreak{} and never used left exactness of F. Thus (2.2)\allowbreak{}–\allowbreak{}(2.4)\allowbreak{},\allowbreak{} convergence,\allowbreak{} and the explicit bounds hold for every additive F with a given Cartan–\allowbreak{}Eilenberg resolution. In that statement R\textasciicircum{}0F(M)\allowbreak{} is not silently identified with F(M)\allowbreak{}.

For a strict example,\allowbreak{} fix a prime ell and let F(M)\allowbreak{}\allowbreak{}=M direct-sum (M/\allowbreak{}ell M)\allowbreak{} on abelian groups. It is additive. On multiplication by ell from Z to Z,\allowbreak{} its second component is zero on Z/\allowbreak{}ell Z,\allowbreak{} so F does not preserve that monomorphism and is not left exact. An injective group J is ell-divisible:\allowbreak{} extend the map ell Z\allowbreak{}->J sending ell to any given y across ell Z\allowbreak{}->Z. The image x of 1 satisfies ell x\allowbreak{}=y. Hence F(J)\allowbreak{}\allowbreak{}=J direct-sum 0. Applied to an injective resolution of M it gives that resolution direct-sum the zero complex. Its derived value has H\textasciicircum{}0\allowbreak{}=M direct-sum 0 and no positive cohomology,\allowbreak{} whereas F(M)\allowbreak{} can have a nonzero second summand. This preserves both summands and shows that the extension is strictly broader than the left-exact source statement.

\subsection{3. Functoriality with the exact page change}\label{proofs-7096-7578-functoriality-with-the-exact-page-change}

Original 7253–\allowbreak{}7262 points forward to the filtered construction. The current forward passage,\allowbreak{} 8835–\allowbreak{}8969,\allowbreak{} was read as a dependency,\allowbreak{} not as completed review of the intervening section. Here are the comparison maps that justify the use made here. The first spectral sequence is canonical starting at E\_\allowbreak{}1; the second is canonical starting at E\_\allowbreak{}2. An arbitrarily chosen injective resolution does not make its raw E\_\allowbreak{}0 terms canonical. The later source\textquotesingle s “r>\allowbreak{}=1” for the filtered construction makes this convention explicit.

Equip K first with S\textasciicircum{}aK\allowbreak{}=sigma\_\allowbreak{}\{>\allowbreak{}=a\}K. Give T the column filtration S\textasciicircum{}aT\allowbreak{}=Tot(I\textasciicircum{}\{p,\allowbreak{}q\}:\allowbreak{}p>\allowbreak{}=a)\allowbreak{}. Its associated graded in degree a is the vertical resolution I\textasciicircum{}\{a,\allowbreak{}\textbackslash{}bullet\}[-a]\allowbreak{},\allowbreak{} with vertical sign (-1)\allowbreak{}\textasciicircum{}a. The augmentation on that graded piece resolves K\textasciicircum{}a[-a]\allowbreak{}. Each T\^{}m has a finite split filtration whose graded objects are injective. Thus K\allowbreak{}->T is a filtered injective resolution for S.

For the second construction put G\textasciicircum{}sK\allowbreak{}=tau\_\allowbreak{}\{<\allowbreak{}=-s\}K,\allowbreak{} with its actual cycles in top degree. Put

\[
 G^sT=\operatorname{Tot}(\tau_{\le -s}^{\,\mathrm{horizontal}} I).
                                                                    \tag{3.1}
\]

This is a subcomplex:\allowbreak{} vertical differentials preserve horizontal cycles and the truncated horizontal differential vanishes on the last cycles. For k\allowbreak{}=-s,\allowbreak{} the associated graded horizontal complex has J\_\allowbreak{}B\textasciicircum{}\{k,\allowbreak{}\textbackslash{}bullet\} in horizontal degree k-1 and J\_\allowbreak{}Z\textasciicircum{}\{k,\allowbreak{}\textbackslash{}bullet\} in horizontal degree k,\allowbreak{} with the actual inclusion b\_\allowbreak{}k between them. Its total differential keeps the signs (-1)\allowbreak{}\textasciicircum{}\{k-1\} and (-1)\allowbreak{}\textasciicircum{}k in these two columns. Each associated graded object of T\^{}m is thus a finite direct sum of injectives. The induced augmentation from gr\_\allowbreak{}G\textasciicircum{}s K resolves its two terms B\^{}k and Z\^{}k by columns. The finite-diagonal comparison of section 2 proves it is a quasi-isomorphism on the total complexes. Moreover the quotient by the first column pair is J\_\allowbreak{}H\textasciicircum{}\{k,\allowbreak{}\textbackslash{}bullet\} in horizontal degree k:\allowbreak{} its kernel is the total complex of the identity on J\_\allowbreak{}B\textasciicircum{}\{k,\allowbreak{}\textbackslash{}bullet\},\allowbreak{} which is contractible. Explicitly,\allowbreak{} on the two copies in horizontal degrees k-1,\allowbreak{}k take the degree -1 map from the second copy to the first to be the identity and zero on the first; the vertical cross terms cancel by (-1)\allowbreak{}\textasciicircum{}\{k-1\}\allowbreak{}+(-1)\allowbreak{}\textasciicircum{}k\allowbreak{}=0. It follows that the graded total is quasi-isomorphic to J\_\allowbreak{}H\textasciicircum{}\{k,\allowbreak{}\textbackslash{}bullet\}[-k]\allowbreak{}. This also proves directly the asserted cohomology description of the corresponding graded source.

For completeness,\allowbreak{} the needed filtered comparison lemma follows from the same extension argument as the injective comparison already proved,\allowbreak{} with strictly exact filtered sequences. An object J with finite filtration and injective graded pieces is a split finite sum of its graded pieces:\allowbreak{} descending induction splits 0\allowbreak{}->F\textasciicircum{}\{a\allowbreak{}+1\}J\allowbreak{}->F\textasciicircum{}aJ\allowbreak{}->gr\textasciicircum{}a J\allowbreak{}->0 because F\textasciicircum{}\{a\allowbreak{}+1\}J is a finite sum of injectives. A filtered map A\allowbreak{}->J from a strict subobject A of B can be extended component by component. The component to the piece in filtration degree a factors through A/\allowbreak{}F\textasciicircum{}\{a\allowbreak{}+1\}A,\allowbreak{} which is a subobject of B/\allowbreak{}F\textasciicircum{}\{a\allowbreak{}+1\}B by strictness; use injectivity of that piece to extend it. The finitely many extensions assemble to a filtered extension.

An acyclic complex on associated graded objects is strictly acyclic with its induced cycle filtrations. To check this,\allowbreak{} lift a filtered cycle through the differential successively on the graded pieces; subtract the resulting boundaries and raise the filtration index at every step. The finite filtration in the degrees involved terminates the process. Its kernels and images are therefore the filtered cycles,\allowbreak{} and its cycle sequences are strict exact sequences. Given a chain map from such a complex to a bounded-below complex of filtered injectives,\allowbreak{} construct a null homotopy upwards from the first nonzero target degree:\allowbreak{} the residual map vanishes on cycles by the previously constructed homotopy equation,\allowbreak{} factors through the next cycles,\allowbreak{} and extends to the next entire term by the filtered extension property. This is the same equation f\textasciicircum{}m\allowbreak{}=d\_\allowbreak{}J\textasciicircum{}\{m-1\}h\textasciicircum{}m\allowbreak{}+h\textasciicircum{}\{m\allowbreak{}+1\}d\textasciicircum{}m proved in the preceding note,\allowbreak{} now with every map filtered. It proves that a filtered quasi-isomorphism induces a bijection on homotopy classes into these targets,\allowbreak{} by applying the Hom sequence to its strictly acyclic cone and its shift.

Consequently a map of source filtered complexes gives a unique homotopy class of comparison maps between their filtered injective resolutions. Identity and composition hold by that uniqueness. An additive F preserves their filtered homotopies:\allowbreak{} all term filtrations of the resolutions split,\allowbreak{} so their images define finite filtrations and F of every filtered component remains filtered. A filtered homotopy induces an ordinary homotopy on associated graded complexes and therefore equal maps on E\_\allowbreak{}1 and all later pages. This establishes choice independence and naturality for both filtered constructions,\allowbreak{} even when a given resolution is available without a global enough-injectives assumption.

Here is the exact relation of (3.1)\allowbreak{} to the second double-complex sequence,\allowbreak{} rather than just an agreement of named E\_\allowbreak{}2 objects. Let V\^{}aT be the vertical-column filtration,\allowbreak{} retaining summands I\textasciicircum{}\{p,\allowbreak{}q\} with q>\allowbreak{}=a. For a filtered complex (C,\allowbreak{}V)\allowbreak{} define

\[
 (\operatorname{Dec}V)^s C^m
   =V^{s+m}C^m\cap d^{-1}(V^{s+m+1}C^{m+1}).                  \tag{3.2}
\]

For C\allowbreak{}=T the first condition retains horizontal degrees p<\allowbreak{}=-s; the second requires that the component in horizontal degree -s be a horizontal cycle. All remaining components already satisfy it,\allowbreak{} including their vertical differentials. Hence Dec V\allowbreak{}=G with the actual subobjects (3.1)\allowbreak{}. After applying F the same equality holds:\allowbreak{} the required horizontal kernels are preserved by the split sequences (2.3)\allowbreak{}.

To verify pages and differentials,\allowbreak{} use \[
 Z_r^{a,m}(V)=V^a C^m\cap d^{-1}V^{a+r}C^{m+1},\qquad
 E_r^{a,m-a}(V)=
 \frac{Z_r^{a,m}(V)}
 {Z_{r-1}^{a+1,m}(V)+dZ_{r-1}^{a-r+1,m-1}(V)}
 \quad(r\ge1).                                               \tag{3.3}
\] Formula (3.2)\allowbreak{} and d\textasciicircum{}2\allowbreak{}=0 give Z\_\allowbreak{}r\textasciicircum{}\{s,\allowbreak{}m\}(Dec V)\allowbreak{}\allowbreak{}=Z\_\allowbreak{}\{r\allowbreak{}+1\}\textasciicircum{}\{s\allowbreak{}+m,\allowbreak{}m\}(V)\allowbreak{}. The two denominator terms become respectively Z\_\allowbreak{}r\textasciicircum{}\{s\allowbreak{}+m\allowbreak{}+1,\allowbreak{}m\}(V)\allowbreak{} and dZ\_\allowbreak{}r\textasciicircum{}\{s\allowbreak{}+m-r,\allowbreak{}m-1\}(V)\allowbreak{}. They are exactly the denominator for E\_\allowbreak{}\{r\allowbreak{}+1\}\textasciicircum{}\{s\allowbreak{}+m,\allowbreak{}-s\}(V)\allowbreak{}. The differential is induced by the same original d\_\allowbreak{}C on both quotients. Thus

\[
 E_r^{s,t}(G)= {}''E_{r+1}^{\,2s+t,-s},\qquad r\ge1,            \tag{3.4}
\]

compatibly with differentials and all maps. Its change of indices sends bidegree (r,\allowbreak{}1-r)\allowbreak{} to (r\allowbreak{}+1,\allowbreak{}-r)\allowbreak{},\allowbreak{} as required. At r\allowbreak{}=1,\allowbreak{} (3.4)\allowbreak{} gives precisely the source\textquotesingle s R\textasciicircum{}\{2s\allowbreak{}+t\}F(H\textasciicircum{}\{-s\}K)\allowbreak{}\allowbreak{}=\{\}''E\_\allowbreak{}2\textasciicircum{}\{2s\allowbreak{}+t,\allowbreak{}-s\}. This proves the promised functoriality from E\_\allowbreak{}1 and E\_\allowbreak{}2,\allowbreak{} respectively,\allowbreak{} without asserting it for arbitrary E\_\allowbreak{}0 choices. It also supplies the actual page-shifting comparison used in current Cohomology 7342–\allowbreak{}7367.

\subsection{\texorpdfstring{4. Composition:\allowbreak{} the evaluation object and the additive extension}{4.  the evaluation object and the additive extension}}\label{proofs-7096-7578-composition-the-evaluation-object-and-the-additive-extension}

Write F:\allowbreak{}A\allowbreak{}->B and G:\allowbreak{}B\allowbreak{}->C and assume A and B have enough injectives. Keep the canonical comparison t:\allowbreak{}R(GF)\allowbreak{}\allowbreak{}->RG RF constructed in section 1 of PROOFS\_\allowbreak{}5091\_\allowbreak{}5900.md. For X represented by A\^{}\textbackslash bullet and a resolution A\textasciicircum{}\textbackslash{}bullet\allowbreak{}->I\textasciicircum{}\textbackslash{}bullet,\allowbreak{} its component is exactly

\[
 t_X:\quad
 R(GF)(X)=G(F(I^\bullet))
       \xrightarrow{\eta_{G,F(I^\bullet)}}RG(F(I^\bullet))
       =RG(RF(X)).                                            \tag{4.1}
\]

The displayed equalities are the specified injective-resolution identifications,\allowbreak{} not an assertion that a nonexact F preserves arbitrary quasi-isomorphisms.

If every F(I)\allowbreak{},\allowbreak{} for I injective in A,\allowbreak{} is right G-acyclic,\allowbreak{} all terms of F(I\textasciicircum{}\textbackslash{}bullet)\allowbreak{} are G-acyclic. The already proved Leray acyclicity lemma makes the middle arrow of (4.1)\allowbreak{} an isomorphism. Conversely evaluate t at the object I[0]\allowbreak{} of D\textasciicircum{}\allowbreak{}+(A)\allowbreak{}. Its two sides identify with G(F(I)\allowbreak{})\allowbreak{}[0]\allowbreak{} and RG(F(I)\allowbreak{}[0]\allowbreak{})\allowbreak{},\allowbreak{} and its component is their canonical comparison. Thus an isomorphism t makes F(I)\allowbreak{} right G-acyclic.

\textbf{DERIVED-NEW-026.} Original 7311 instead says to evaluate t on RF(I)\allowbreak{}\allowbreak{}=F(I)\allowbreak{}. That is an object of D\textasciicircum{}\allowbreak{}+(B)\allowbreak{},\allowbreak{} whereas t has domain D\textasciicircum{}\allowbreak{}+(A)\allowbreak{}; for unrelated A and B this is not an allowable argument of t. Replace that clause by evaluation at I[0]\allowbreak{},\allowbreak{} using RF(I[0]\allowbreak{})\allowbreak{}\allowbreak{}=F(I)\allowbreak{}[0]\allowbreak{}. The argument above proves the conclusion with exactly that component. The theorem is valid; it is the evaluation instruction in the proof that is ill-typed.

\textbf{DERIVED-UNDER-011,\allowbreak{} composition part.} This equivalence also holds for additive F and G without left exactness. Enough injectives defines all three right derived functors; (4.1)\allowbreak{} and Leray acyclicity use additivity and the actual right-acyclic comparison. For this extension ``right G-acyclic'' means that G(M)\allowbreak{}[0]\allowbreak{}\allowbreak{}->RG(M[0]\allowbreak{})\allowbreak{} itself is an isomorphism. Vanishing R\textasciicircum{}iG(M)\allowbreak{} for i\textgreater0 alone is insufficient:\allowbreak{} it omits the degree-zero comparison,\allowbreak{} as the example in section 2 illustrates.

Under these equivalent conditions choose a Cartan–\allowbreak{}Eilenberg resolution of F(I\textasciicircum{}\textbackslash{}bullet)\allowbreak{},\allowbreak{} and apply section 2 with functor G. Its second sequence has \[
 E_2^{p,q}=R^pG(H^q(F(I^\bullet)))=R^pG(H^q(RF(X))).
                                                               \tag{4.2}
\] Its abutment is H\textasciicircum{}\{p\allowbreak{}+q\}(RG(RF(X)\allowbreak{})\allowbreak{})\allowbreak{},\allowbreak{} identified by H(t\_\allowbreak{}X)\allowbreak{} with H\textasciicircum{}\{p\allowbreak{}+q\}(R(GF)\allowbreak{}(X)\allowbreak{})\allowbreak{}. This proves the stated Grothendieck sequence with the exact comparison map. If I\^{}\textbackslash bullet is zero below n,\allowbreak{} F(I\textasciicircum{}\textbackslash{}bullet)\allowbreak{} is zero below n; take the Cartan–\allowbreak{}Eilenberg resolution with the same horizontal bound. Section 2 supplies page m-n\allowbreak{}+2 stabilization and at most m-n\allowbreak{}+1 graded pieces in each degree m>\allowbreak{}=n. Naturality from E\_\allowbreak{}2 follows from section 3 and (4.1)\allowbreak{}.

The report dispositions in this interval are specific:\allowbreak{} RECON-045 /\allowbreak{} P02-E0050 /\allowbreak{} DERIVED-041 is existing MC-STK-ERR-0852 (``be'' to ``are'' at 7278)\allowbreak{}. RECON-046 /\allowbreak{} P02-E0051 /\allowbreak{} DERIVED-042 is existing MC-STK-ERR-0853,\allowbreak{} both the premature period and missing article. RECON-047 /\allowbreak{} P02-E0052 /\allowbreak{} DERIVED-043 is existing MC-STK-ERR-0854,\allowbreak{} the extra parenthesis in (4.1)\allowbreak{}. RECON-048 /\allowbreak{} P02-E0053 /\allowbreak{} DERIVED-044 is existing MC-STK-ERR-0855. P02 proposes joining the fragment to the next sentence,\allowbreak{} whereas the existing two-operation edit makes it a complete assumption sentence. Both address the same defect; the existing edit retains exactly the original mathematical hypotheses and is not replaced solely to match P02\textquotesingle s wording.

\subsection{5. Finite-prefix propagation of DERIVED-UNDER-007}\label{proofs-7096-7578-finite-prefix-propagation-of-derived-under-007}

The stronger finite-prefix comparison already proved in PART07 has a precise receiver in (4.1)\allowbreak{}. Let F and G be additive,\allowbreak{} A and B have enough injectives,\allowbreak{} and fix X\allowbreak{}->I\textasciicircum{}\textbackslash{}bullet. It suffices here that the particular F(I\textasciicircum{}j)\allowbreak{} are right G-acyclic for j<\allowbreak{}=m; no assertion about all injectives is needed. Applying DERIVED-UNDER-007 to the actual complex F(I\textasciicircum{}\textbackslash{}bullet)\allowbreak{} and functor G gives \[
 H^j(t_X)\text{ is an isomorphism for }j<m,\qquad
 H^m(t_X)\text{ is monic}.                                    \tag{5.1}
\] Its boundary defect is exactly the intersection from the previous proof:\allowbreak{} with A\allowbreak{}=F(I\textasciicircum{}\textbackslash{}bullet)\allowbreak{},\allowbreak{} P\allowbreak{}=sigma\_\allowbreak{}\{<\allowbreak{}=m\}A and T\allowbreak{}=sigma\_\allowbreak{}\{>\allowbreak{}=m\allowbreak{}+1\}A,\allowbreak{} the cokernel of H\textasciicircum{}m(t\_\allowbreak{}X)\allowbreak{} is identified with im(delta\_\allowbreak{}G)\allowbreak{} intersect ker H\textasciicircum{}\{m\allowbreak{}+1\}(eta\_\allowbreak{}\{G,\allowbreak{}T\})\allowbreak{},\allowbreak{} where delta\_\allowbreak{}G:\allowbreak{}H\textasciicircum{}m(G(P)\allowbreak{})\allowbreak{}\allowbreak{}->H\textasciicircum{}\{m\allowbreak{}+1\}(G(T)\allowbreak{})\allowbreak{} is the boundary of the original termwise split triangle T\allowbreak{}->A\allowbreak{}->P\allowbreak{}->T[1]\allowbreak{}. This follows by the same comparison of its two exact cohomology sequences; every object and map is now substituted from (4.1)\allowbreak{}.

If G is also left exact,\allowbreak{} the complete PART07 addendum improves (5.1)\allowbreak{} to \[
 H^j(t_X)\text{ is an isomorphism for }j\le m+1,\qquad
 H^{m+2}(t_X)\text{ is monic}.                                \tag{5.2}
\] For clarity,\allowbreak{} the extra degrees come from the bottom of T:\allowbreak{} left exactness makes H\textasciicircum{}\{m\allowbreak{}+1\}(G(T)\allowbreak{})\allowbreak{}\allowbreak{}->H\textasciicircum{}\{m\allowbreak{}+1\}(RG(T)\allowbreak{})\allowbreak{} an isomorphism and the next comparison monic. Comparing cokernels of H\textasciicircum{}m(G(P)\allowbreak{})\allowbreak{}\allowbreak{}->H\textasciicircum{}\{m\allowbreak{}+1\}(G(T)\allowbreak{})\allowbreak{} proves the degree m\allowbreak{}+1 assertion,\allowbreak{} since P computes RG(P)\allowbreak{}. The adjacent cohomology objects of P vanish in the next degree,\allowbreak{} proving the last assertion. This retains all of A,\allowbreak{} including its terms above m. It is not a claim that the truncated term complex alone computes every displayed cohomology object.

The Grothendieck sequence without the all-terms acyclicity hypothesis still abuts to RG(RF(X)\allowbreak{})\allowbreak{}; (5.1)\allowbreak{} or (5.2)\allowbreak{} identifies only the proved degree range with R(GF)\allowbreak{}(X)\allowbreak{}. No global abutment identification is asserted from a finite prefix.

\subsection{6. The equivalence and the meaning of uniqueness}\label{proofs-7096-7578-the-equivalence-and-the-meaning-of-uniqueness}

The full subcategory of injectives is additive and closed under finite sums. Shifts and mapping cones of bounded-below complexes in it stay there. The canonical functor q:\allowbreak{}K\textasciicircum{}\allowbreak{}+(I)\allowbreak{}\allowbreak{}->D\textasciicircum{}\allowbreak{}+(A)\allowbreak{} therefore preserves distinguished cone triangles and shifts. Every object of D\textasciicircum{}\allowbreak{}+(A)\allowbreak{} has a bounded below representative and an injective resolution,\allowbreak{} giving essential surjectivity. For complexes U,\allowbreak{}V in K\textasciicircum{}\allowbreak{}+(I)\allowbreak{},\allowbreak{} the injective Hom comparison proves Hom\_\allowbreak{}\{K(A)\allowbreak{}\}(U,\allowbreak{}V)\allowbreak{}\allowbreak{}=Hom\_\allowbreak{}\{D(A)\allowbreak{}\}(U,\allowbreak{}V)\allowbreak{}. The bounded homotopy and derived subcategories are full on these objects,\allowbreak{} so this is the required full faithfulness. This proves original 7377–\allowbreak{}7397 with its actual functor.

RECON-049 /\allowbreak{} DERIVED-046 is the existing MC-STK-ERR-0857 ``Full faithfulness'' correction at 7396. This also closes the second locus of P02-E0039 /\allowbreak{} OCC-01511. Its first locus 6197–\allowbreak{}6198 was proved in section 4 of PROOFS\_\allowbreak{}5910\_\allowbreak{}6199.md and linked to RECON-035 /\allowbreak{} MC-STK-ERR-0840. The physical P02 report is counted once,\allowbreak{} in RECON-049,\allowbreak{} with an explicit link to that earlier component. The two French reports remain their distinct physical records.

A resolution functor was defined at 7407–\allowbreak{}7418 as the pair of choices (jK,\allowbreak{}i\_\allowbreak{}K:\allowbreak{}K\allowbreak{}->jK)\allowbreak{},\allowbreak{} not as j without augmentation. Precomposition with i\_\allowbreak{}K gives a bijection \[
 \operatorname{Hom}_{K(A)}(jK,jL)
       \longrightarrow\operatorname{Hom}_{K(A)}(K,jL).
                                                               \tag{6.1}
\] Let j(alpha)\allowbreak{} be the unique preimage of i\_\allowbreak{}L alpha. For the identity it is the identity; for beta alpha it is j(beta)\allowbreak{}j(alpha)\allowbreak{},\allowbreak{} since that composite has the required precomposition. Additivity follows from the same uniqueness and the additive map in (6.1)\allowbreak{}. Thus j is an additive functor and i a natural transformation in K\textasciicircum{}\allowbreak{}+(A)\allowbreak{}. Applying localization gives a natural isomorphism since each i\_\allowbreak{}K is a quasi-isomorphism.

For another pair (j',\allowbreak{}i')\allowbreak{},\allowbreak{} the same bijection gives a unique a\_\allowbreak{}K:\allowbreak{}jK\allowbreak{}->j'K with a\_\allowbreak{}K i\_\allowbreak{}K\allowbreak{}=i'\_\allowbreak{}K in K\textasciicircum{}\allowbreak{}+(A)\allowbreak{}. Construct b\_\allowbreak{}K in the opposite direction. The composites b\_\allowbreak{}Ka\_\allowbreak{}K and a\_\allowbreak{}Kb\_\allowbreak{}K are identities by (6.1)\allowbreak{}. Naturality is proved by precomposing the two candidate maps jK\allowbreak{}->j'L with i\_\allowbreak{}K. This proves existence and uniqueness of the augmentation-compatible natural isomorphism a. The existence statement for a set of objects follows by choosing one injective resolution for each such object.

The unqualified phrase at 7463–\allowbreak{}7464 is read using the defined structured object (j,\allowbreak{}i)\allowbreak{}. Without compatibility with i,\allowbreak{} uniqueness among all natural isomorphisms would be false. For example,\allowbreak{} over Q-vector spaces every object is injective and the identity functor on K\textasciicircum{}\allowbreak{}+ has distinct natural automorphisms 1 and -1; they differ on Q[0]\allowbreak{}. Only 1 is compatible with the identity augmentation. The source\textquotesingle s definition and its explicit comparison at 7571–\allowbreak{}7576 already provide this compatibility,\allowbreak{} so no additional source error is counted. This records the adverse interpretation and its exact counterexample without turning shorthand into an unsupported new defect.

RECON-050 /\allowbreak{} DERIVED-047 is existing MC-STK-ERR-0858,\allowbreak{} ``resolution functors'' at 7570. RECON-051 /\allowbreak{} P02-E0054 /\allowbreak{} DERIVED-045 is existing MC-STK-ERR-0856:\allowbreak{} with i:\allowbreak{}K\allowbreak{}->I,\allowbreak{} i':\allowbreak{}K\allowbreak{}->J and a:\allowbreak{}I\allowbreak{}->J,\allowbreak{} the equality is i'\allowbreak{}=a i in K\textasciicircum{}\allowbreak{}+(A)\allowbreak{}. The arrow a lies in K\textasciicircum{}\allowbreak{}+(I)\allowbreak{},\allowbreak{} but K need not. Both the composition-order and ambient-category repairs are needed and retained,\allowbreak{} even though P02 reports only the first.

\subsection{\texorpdfstring{7. Exactness,\allowbreak{} shifts,\allowbreak{} and the quasi-inverse}{7.   and the quasi-inverse}}\label{proofs-7096-7578-exactness-shifts-and-the-quasi-inverse}

There is a unique xi\_\allowbreak{}K:\allowbreak{}j(K[1]\allowbreak{})\allowbreak{}\allowbreak{}->j(K)\allowbreak{}[1]\allowbreak{} with \[
 \xi_K i_{K[1]}=i_K[1].
                                                               \tag{7.1}
\] Both arrows on the right and on the left\textquotesingle s domain are resolutions of K[1]\allowbreak{}. The comparison in section 6 makes xi\_\allowbreak{}K invertible. For alpha:\allowbreak{}K\allowbreak{}->L,\allowbreak{} both xi\_\allowbreak{}L j(alpha[1]\allowbreak{})\allowbreak{} and j(alpha)\allowbreak{}[1]\allowbreak{}xi\_\allowbreak{}K become i\_\allowbreak{}L[1]\allowbreak{}alpha[1]\allowbreak{} after precomposition with i\_\allowbreak{}\{K[1]\allowbreak{}\}; the same uniqueness proves their equality. This proves naturality,\allowbreak{} including the source\textquotesingle s shift convention d\_\allowbreak{}\{K[1]\allowbreak{}\}\textasciicircum{}m\allowbreak{}=-d\_\allowbreak{}K\textasciicircum{}\{m\allowbreak{}+1\}.

For a distinguished triangle (K,\allowbreak{}L,\allowbreak{}M,\allowbreak{}f,\allowbreak{}g,\allowbreak{}h)\allowbreak{},\allowbreak{} the triangle of chosen resolutions has final arrow xi\_\allowbreak{}K j(h)\allowbreak{}. Its image under q is isomorphic to the original localized triangle via the three i\textquotesingle s and (7.1)\allowbreak{}. It is consequently distinguished in D\textasciicircum{}\allowbreak{}+(A)\allowbreak{}. To see directly that q reflects this property,\allowbreak{} complete the first arrow in K\textasciicircum{}\allowbreak{}+(I)\allowbreak{} to a cone triangle. In D\textasciicircum{}\allowbreak{}+(A)\allowbreak{} the cone triangle and the proposed triangle have identical first arrows,\allowbreak{} and the triangle axiom gives a map on their third objects. It is an isomorphism:\allowbreak{} applying each representable cohomological functor to the two long exact sequences gives this by the five lemma. Full faithfulness of q lifts that map and its inverse to K\textasciicircum{}\allowbreak{}+(I)\allowbreak{},\allowbreak{} and lifts its commutative triangle equalities. Thus the proposed triangle is isomorphic to a cone triangle in K\textasciicircum{}\allowbreak{}+(I)\allowbreak{},\allowbreak{} proving that (j,\allowbreak{}xi)\allowbreak{} is exact.

\begingroup\raggedright
For a quasi-isomorphism alpha,\allowbreak{} naturality of i shows j(alpha)\allowbreak{}i\_\allowbreak{}K\allowbreak{}=i\_\allowbreak{}L alpha is a quasi-isomorphism. Since i\_\allowbreak{}K is one too,\allowbreak{} j(alpha)\allowbreak{} is a quasi-isomorphism. Between injective complexes it is invertible in K\textasciicircum{}\allowbreak{}+(I)\allowbreak{}:\allowbreak{} its image in D\textasciicircum{}\allowbreak{}+(A)\allowbreak{} is invertible and q is fully faithful. The localization universal property therefore factors j through Q:\allowbreak{}K\textasciicircum{}\allowbreak{}+(A)\allowbreak{}\allowbreak{}->D\textasciicircum{}\allowbreak{}+(A)\allowbreak{} as j'Q,\allowbreak{} uniquely with this specified equality. For U in K\textasciicircum{}\allowbreak{}+(I)\allowbreak{},\allowbreak{} i\_\allowbreak{}U:\allowbreak{}U\allowbreak{}->jU is itself a morphism in K\textasciicircum{}\allowbreak{}+(I)\allowbreak{} and is invertible; these maps give Id\allowbreak{}->j'q. For a representative K,\allowbreak{} Q(i\_\allowbreak{}K)\allowbreak{} gives QK\allowbreak{}->qjK. This transformation descends through Q:\allowbreak{} naturality holds on original maps by section 6 and on inverted quasi-isomorphisms by inverting that same naturality equation. It gives Id\allowbreak{}->qj' on D\textasciicircum{}\allowbreak{}+(A)\allowbreak{}. Both are natural isomorphisms,\allowbreak{} proving the quasi-inverse assertion omitted at 7560–\allowbreak{}7561.
\par\endgroup

The reference to a square in 7554–\allowbreak{}7556 uses the naturality square just proved; although its explicit alpha-labelled display occurs in the earlier resolution-functor lemma,\allowbreak{} that relation is also used in the exactness proof. This does not require a new mathematical correction. For a category with a proper class of objects the comparison proof still works for given choices; set-sized choice alone does not supply a class of such choices. The source explicitly defers that existence issue to its next section. This review does not assert that the issue has already been resolved.

\subsection{8. Receiving statements and propagation limits}\label{proofs-7096-7578-receiving-statements-and-propagation-limits}

Current Cohomology 2140–\allowbreak{}2193 applies composition to the actual pair f\_\allowbreak{}* and Gamma(Y,\allowbreak{}-)\allowbreak{},\allowbreak{} retaining the restriction-of-scalars functor along O\_\allowbreak{}Y(Y)\allowbreak{}\allowbreak{}->O\_\allowbreak{}X(X)\allowbreak{}. On a chosen injective resolution I\^{}\textbackslash bullet the comparison is \[
 \Gamma(Y,f_*I^\bullet)
     \longrightarrow R\Gamma(Y,f_*I^\bullet).
                                                               \tag{8.1}
\] Its left side is Gamma(X,\allowbreak{}I\textasciicircum{}\textbackslash{}bullet)\allowbreak{} with that specified scalar restriction. Section 4 evaluates acyclicity on I[0]\allowbreak{} in the category of O\_\allowbreak{}X-modules; it does not try to put f\_\allowbreak{}*I into that category. The receiver\textquotesingle s injective acyclicity hypothesis is exactly the one used in (4.1)\allowbreak{}. Its Leray sequence at 2261–\allowbreak{}2283 therefore retains E\_\allowbreak{}2\textasciicircum{}\{p,\allowbreak{}q\}\allowbreak{}=H\textasciicircum{}p(Y,\allowbreak{}R\textasciicircum{}qf\_\allowbreak{}*\textbackslash{}mathcal F\textasciicircum{}\textbackslash{}bullet)\allowbreak{} and its stated abutment via (8.1)\allowbreak{}. The correction of 7248–\allowbreak{}7249 changes no E\_\allowbreak{}2 statement here:\allowbreak{} it repairs the resolving complex in the proof.

For a chosen I\^{}\textbackslash bullet zero below n the spectral sequence has at most m-n\allowbreak{}+1 graded pieces in degree m and stabilizes there by page m-n\allowbreak{}+2. If only the prefix f\_\allowbreak{}*I\textasciicircum{}j,\allowbreak{} j<\allowbreak{}=m,\allowbreak{} is known Gamma(Y,\allowbreak{}-)\allowbreak{}-acyclic,\allowbreak{} (5.2)\allowbreak{} still proves that (8.1)\allowbreak{} is an isomorphism on cohomology through degree m\allowbreak{}+1 and monic in degree m\allowbreak{}+2. The existing receiver supplies acyclicity of every term,\allowbreak{} so this is an editorial strengthening,\allowbreak{} not a change to its hypotheses or statement.

Current Sites Cohomology 2304–\allowbreak{}2344 has the same exact receiver with Gamma(D,\allowbreak{}-)\allowbreak{},\allowbreak{} f\_\allowbreak{}* and Gamma(C,\allowbreak{}-)\allowbreak{}. Its cited total acyclicity of f\_\allowbreak{}*I supplies the full right-acyclic condition. Substitution into (4.1)\allowbreak{},\allowbreak{} (4.2)\allowbreak{},\allowbreak{} and (5.2)\allowbreak{} proves the identical comparison,\allowbreak{} page bounds,\allowbreak{} and finite-prefix refinement for this receiver,\allowbreak{} retaining its ringed topoi and module categories.

Current Cohomology 7342–\allowbreak{}7367 explicitly uses the reindexing p\allowbreak{}=-j,\allowbreak{}q\allowbreak{}=i\allowbreak{}+2j for the cohomology-of-cohomology sequence. Formula (3.4)\allowbreak{} supplies the page change and actual differential as well as that reindexing. Its first sequence at 7370–\allowbreak{}7391 uses the filtration S,\allowbreak{} whose graded injective resolution was checked in section 3. Thus the two identifications are compatible with maps from their canonical pages. The bounded-below cases here are proved; this note does not infer unbounded convergence from those bounds.

A root-chapter reference search also found other users of the two source lemmas. Those hits are dependency routes,\allowbreak{} not claims that their whole proofs were read. The source statements of the corrected lemmas are unchanged. The precise receiver calculations above and the forward filtered comparison are the propagation completed in this part. Joint research synthesis remains later than the correction and received-theorem queue.
